% Criterion 9 converted from the supplied Word source.
% Expected packages: array, booktabs, longtable, graphicx, hyperref, calc, enumitem, float.
\providecommand{\real}[1]{#1}
\providecommand{\tightlist}{}

\hypertarget{criterion-9-industrial-linkages}{%
\chapter{Criterion 9: Industrial Linkages}\label{criterion-9-industrial-linkages}}

NUST recognises industry--academia collaboration as an important
mechanism for aligning academic programs and research with
industrial needs, promoting innovation, and facilitating the
application of research outcomes. Such collaboration provides
opportunities for industry to engage with university researchers,
access specialised expertise and talent, and explore solutions to
practical challenges. It also enables faculty and students to gain
exposure to industrial practices, emerging technologies, and market
requirements.

NUST has established institutional mechanisms to facilitate
research collaboration, industry engagement, intellectual property
management, technology transfer, entrepreneurship, and graduate
employability. Key mechanisms include:

\begin{enumerate}
    \item \textbf{Corporate Advisory Council (CAC):}
    Facilitates engagement between NUST and industry through
    corporate partnerships, industry consultation, collaborative
    research, and activities connecting students and faculty with
    industrial organisations.

    \item \textbf{Innovation \& Commercialisation Office NUST (ICON):}
    Provides an institutional framework for connecting research
    and innovation with industrial application and commercial
    opportunities. Its constituent offices include the Corporate
    Advisory Council (CAC), NUST Intellectual Property Office
    (NIPO), Technology Transfer Office (TTO), and Business
    Development Office (BDO).

    \item \textbf{Directorate for Research:}
    Supports research-related activities and contributes to the
    development of NUST's research capabilities and institutional
    research agenda.

    \item \textbf{National Science \& Technology Park (NSTP):}
    Forms part of NUST's research and innovation ecosystem,
    supporting engagement among researchers, enterprises, and
    technology-oriented organisations.

    \item \textbf{NUST Placement Office and relevant school-level
    liaison arrangements:}
    Support engagement with employers and industry, including
    opportunities relating to student placements, internships,
    industrial exposure, and graduate employment, as applicable
    to the respective unit's mandate.
\end{enumerate}

Through these mechanisms, NUST facilitates interaction between
academic departments, researchers, students, and industry partners.
Activities such as industry consultations, collaborative research,
consultancy, final-year design projects, internships, industrial
visits, intellectual property protection, and technology licensing
provide pathways for translating academic expertise into practical
applications. The nature and outcomes of these collaborations
depend on the requirements of individual industry partners and
the scope of each engagement. Evidence of these linkages is reflected in documented industry
collaborations, research and consultancy activities, student
engagement with industry, and technology-transfer initiatives.
Their contribution to innovation, commercialisation, and graduate
employability should be demonstrated through relevant institutional
records and examples.

\section{Industrial Collaborations at NUST-SEECS}

At the school level, the School of Electrical Engineering and
Computer Science (SEECS) engages with organisations from the public
and private sectors, as well as national and international research
and technology partners. These linkages provide opportunities for
industry-oriented research, consultancy, applied technology
development, and student engagement with real-world problems.
Some of the organisations associated with SEECS collaborations, projects and professional engagement include the following:

\begin{enumerate}[noitemsep, topsep=0pt]
    \item Microsoft
    \item Wacom Ltd., Japan
    \item Dynasys
    \item Fauji Fertilizers Company
    \item Cisco
    \item IBM
    \item NESCOM
    \item Teradata
    \item SkyElectric Pakistan
    \item German Research Center for Artificial Intelligence (DFKI),
          Germany
    \item Huawei
    \item Acrux Technologies
    \item National Agricultural Research Centre (NARC)
    \item Walee Technologies
    \item One Screen
\end{enumerate}

SEECS maintains industry linkages to support the relevance, quality,
and continuous improvement of the BE Software Engineering (BESE)
program in accordance with applicable Pakistan Engineering Council
(PEC) accreditation requirements. At the program level, engagement
with industry is supported through the Industrial Advisory Board
(IAB), collection and analysis of employer and industry feedback,
and consultation on curriculum relevance and graduate competencies.
These mechanisms contribute to the review of Program Educational
Objectives (PEOs) and help identify emerging professional skills
and changing workplace requirements.

Industry engagement is also reflected in student internships,
industry-oriented Final Year Design Projects (FYDPs), project
supervision and mentoring by industry professionals, faculty
consultancy, and applied research addressing practical challenges.
Where supported by documented agreements and project records,
collaborative activities provide opportunities for students and
faculty to work on

```latex
\section{Industry Liaison Office}

The Industrial Liaison Office (ILO) facilitates engagement between
NUST, its constituent schools and colleges, and industry partners.
Its activities support graduate employability, student internships,
industry interaction, and the identification of opportunities for
collaborative research and the application of academic expertise to
industrial challenges. At SEECS, the ILO works with academic
departments and relevant institutional offices to facilitate
interaction between employers, faculty members, and students. The ILO's industry-engagement activities contribute to identifying
employer requirements, connecting organisations with relevant
academic expertise, and facilitating opportunities for students
and graduates. Where applicable, coordination with the Technology
Transfer Office (TTO) and other relevant institutional units can
support the assessment and commercialisation of research outcomes.
The specific responsibilities of the ILO are distinguished from
those of the Industry Advisory Board (IAB), which provides a
program-level mechanism for industry consultation, and the
Career Development Office (CDO), which supports career development
and related student activities.

\subsection{Objectives of the ILO}

The principal objectives of the ILO relevant to SEECS include:

\begin{itemize}[noitemsep, topsep=0pt]
    \item Develop and maintain relationships with industry partners
    and identify opportunities for collaboration relevant to
    SEECS's research and academic activities.

    \item Facilitate the identification of industrial challenges
    and connect relevant organisations with appropriate SEECS departments, research groups, and labs.

    \item Coordinate with employers regarding graduate recruitment,
    placement opportunities, and engagement with prospective
    candidates.

    \item Facilitate student internships by coordinating with
    industry partners and the relevant academic departments.

    \item Support employer recruitment activities, including
    recruitment drives, company presentations, and employer
    interaction sessions, as arranged.

    \item Facilitate the collection of employer feedback on
    graduate competencies and communicate relevant findings to
    the appropriate institutional and academic bodies.

    \item Support alumni engagement and related activities in
    coordination with the relevant NUST offices.

    \item Facilitate industry engagement activities and exhibitions
    in coordination with the relevant institutional units.
\end{itemize}

\subsection{ILO Activities at SEECS}

The ILO facilitates employer engagement and provides opportunities
for organisations to interact with SEECS students and graduates.
Its activities complement the academic department's responsibility
for curriculum delivery, program assessment, and continuous
quality improvement.

\textbf{Employer Sessions and Recruitment Drives:}
Employer sessions provide opportunities for organisations to
introduce career opportunities and interact with prospective
candidates. Depending on the requirements of the recruiting
organisation, these activities may include company presentations,
written tests, interviews, and other recruitment-related
interactions. Such engagement helps connect students and graduates
with prospective employers.

\textbf{Industry Talk Series:}
Industry professionals may be invited to interact with students
through talks and professional engagement sessions. These
activities expose students to industrial practices, emerging
technologies, workplace expectations, and career pathways.
They also provide opportunities for students to understand the
practical application of concepts covered in their academic
programs.

\textbf{Student Development Workshops:}
Professional development activities help students strengthen
communication, teamwork, interpersonal interaction, and other
transferable skills relevant to employment. Workshops and related
sessions may be organised in coordination with the relevant
institutional offices and academic departments, according to
student needs and available opportunities.

\textbf{Internships and Graduate Placement:}
The ILO facilitates interaction with employers regarding internship
and graduate placement opportunities. The relevant academic
departments coordinate program-related requirements and student
participation, while internship experience and employer feedback
can provide useful input for program review and improvement.

\textbf{Industry Feedback and Academic Engagement:}
Feedback from employers and industry representatives contributes
to understanding changing professional requirements and graduate
competencies. At the BE Software Engineering (BESE) program
level, such feedback is considered alongside input from the
Industry Advisory Board (IAB), alumni, and other stakeholders.
The resulting findings can inform the review of program
Educational Objectives (PEOs), curriculum relevance, and
continuous quality improvement. Through these activities, the ILO supports the broader
industry--academia engagement framework at SEECS. Its activities
complement industry consultation through the IAB, collaborative
research, faculty consultancy, and industry engagement in Final
Year Design Projects (FYDPs). 

```latex
\section{Industry Advisory Board}

The Faculty of Computing (FoC) at the School of Electrical Engineering and Computer Science (SEECS), NUST, engages industry representatives through its Industry Advisory Board (IAB) to strengthen the relevance of its academic programs to professional practice and evolving industry requirements. The IAB provides a forum for consultation between industry and faculty representatives on curriculum development, teaching and learning, student projects, graduate attributes, and emerging areas of computing. Its feedback supports continuous improvement across the computing programs and helps align academic outcomes with the expectations of employers and the computing profession.

The FoC IAB comprises invited industry representatives and academic representatives from the computing departments at SEECS. Academic representation includes the Associate Dean (FoC), Heads of the relevant departments, program and Final Year Design Project (FYDP) coordinators, the Industry Liaison Officer, and quality assurance representatives, as applicable to the meeting. Representation from each computing department enables the board to consider both shared industry needs across computing disciplines and program-specific requirements. Industry participation brings perspectives from organisations of different sizes and sectors, including established employers and small and medium-sized enterprises.

The SEECS Industry Liaison Officer serves as secretary of the IAB. The secretary coordinates with the Chair and relevant faculty representatives to prepare meeting agendas, records and circulates the minutes, and follows up on matters arising from the discussions. Recommendations relevant to academic programs, quality assurance, and professional accreditation are documented and communicated to the appropriate departments and faculty-level bodies for consideration. The IAB provides advice and recommendations, formal academic decisions, and accreditation submissions remain the responsibility of the relevant authorised bodies.

The IAB meets at least once annually, with additional meetings convened when required by the Chair or SEECS administration. Meetings may be held in person or online, depending on circumstances. Meeting records document the participants, matters discussed, recommendations, and any agreed follow-up actions.

The FoC IAB supports continuous improvement across SEECS computing programs by providing industry feedback on curriculum design, development, delivery, and review. Recent discussions have covered elective offerings, course credit hours and content, faculty development for industry alignment, student mentorship, access to cloud resources, FYDP-related innovation, and alumni and employer survey instruments. The most recent meeting listed in the records below was held at SEECS on 5 May 2026. Table~\ref{tab:iab_meetings} summarises the IAB meetings held since 2024. Detailed minutes are provided in Annexure~\ref{app:iab_minutes}.

\begin{table}[htbp]
\centering
\caption{FoC-SEECS, IAB Meetings Held Since 2024}
\label{tab:iab_meetings}
\begin{tabular}{|c|l|p{0.58\textwidth}|}
\hline
\textbf{S. No.} & \textbf{Date} & \textbf{Key Focus} \\
\hline
1 & 2 May 2024 & Review of curriculum changes in response to earlier IAB feedback, including the VR/AR elective, increased credit hours for Fundamentals of Computer Programming, and FYDP patenting. \\
\hline
2 & 28 February 2025 & Discussion of proposed electives, including Introduction to Quantum Computing, certification-oriented courses, open-book examinations, and cross-departmental minors. \\
\hline
3 & 7 May 2025 & Discussion of faculty development for industry alignment, course-content distribution across Applied AI, DevOps, and Machine Learning, and student mentorship and access to cloud resources. \\
\hline
4 & 5 May 2026 & Review of the revised alumni and employer survey instruments for program-level feedback and assessment of program Educational Objectives (PEOs). \\
\hline
\end{tabular}
\end{table}

\section{Faculty Placement in Industry}\label{sec:c9-faculty-placement}

To enhance university-industry linkages in advancing innovation and industrial research, NUST Faculty Members are placed in the relevant industry annually, during summer break, under the overall ambit of Faculty Placement Program. Through this program, faculty members are given an opportunity to gain practical insight and industrial exposure leading to enhanced university industry linkages. Purpose of the faculty placement program is:

\begin{itemize} [noitemsep, topsep=0pt]
\item To enhance the industrial exposure of Faculty members by providing practical insight that can be transferred further to students.

\item Study the workflow, industrial practices and standards of the host industry and analyse their problems including both technical and nontechnical issues.

\item Explore potential research and development areas of mutual interest and secure R\&D projects/partnerships for NUST.

\item Secure Jobs, Internships and Trainings for NUST students.

\item Introduce NUST and its capacities to the industries in Pakistan.

\item Promote NUST’s agenda for university-industry linkage and technology transfer.

\item Propose the use of IEDC's platform for conducting training sessions for industrial experts.

\item Build strong ties with the industry and recommend annual continuation of the faculty placement program in same or other relevant industries.
\end{itemize}

Table~\ref{tab:place} lists recent industrial placements of SE faculty members. A sample post-placement report is attached as Annexure~\ref{app:post}.

\begin{longtable}{>{\raggedright\arraybackslash}p{0.07\textwidth} >{\raggedright\arraybackslash}p{0.40\textwidth} >{\raggedright\arraybackslash}p{0.25\textwidth} >{\raggedright\arraybackslash}p{0.20\textwidth}}
\caption{Faculty Placement}\\
\label{tab:place}\\
\toprule
\textbf{S No} & \textbf{Faculty Name} & \textbf{Industry} & \textbf{Year} \\
\midrule
\endfirsthead
\toprule
\textbf{S No} & \textbf{Faculty Name} & \textbf{Industry} & \textbf{Duration} \\
\midrule
\endhead
\midrule
\multicolumn{4}{r}{\emph{Continued on next page}}\\
\endfoot
\bottomrule
\endlastfoot
1 & Dr. Momina Moetesum & vResolv.io  & July 2025 \\
2 & Ms. Naema Asif & Devsinc & July 2026\\
\end{longtable}

\section{Industrial Feedback and PEO Attainment Analysis}
\label{sec:c9-industrial-feedback}

The BE Software Engineering program at SEECS obtains stakeholder feedback to assess the continuing relevance of its Program Educational Objectives (PEOs) and identify opportunities for program improvement. The principal sources of stakeholder input include the Industry Advisory Board (IAB), alumni surveys, employer surveys, and departmental reviews. These sources provide feedback on graduates' software engineering knowledge and skills, teamwork and collaboration, innovation, and professional practice with due consideration for societal, ethical, and environmental aspects. The feedback is considered alongside the program's educational objectives, stakeholder expectations, and applicable accreditation requirements.

The LQEC, under the supervision of the HoD-SE, coordinates the program's assessment and continuous quality improvement activities. Alumni and employer survey data are collected annually, while PEO attainment is evaluated on a four-year cycle. The PEO statements are also subject to periodic review against assessment evidence, stakeholder feedback, and the vision and mission of NUST, SEECS, and the Software Engineering Department.

As part of the 2026 assessment initiative, the alumni and employer survey instruments were redesigned to improve the alignment between survey items and the PEOs. Each revised instrument contains 12 Likert-scale items, with four items mapped to each PEO. The revised instruments followed the documented review sequence of DBS on 20 April 2026, FBS on 21 April 2026, and IAB endorsement on 5 May 2026. During the 2026 cycle, 79 alumni responses and 50 employer responses were received.  

The survey findings, IAB recommendations, and departmental reviews provide inputs to the program's continuous quality improvement (CQI) process. Where an issue requiring improvement is identified, the LQEC and relevant program stakeholders review the available evidence, consider possible causes, and recommend appropriate corrective actions. Depending on the issue, these may include reviewing course content, teaching and assessment practices, CLO--PLO mappings, laboratory or project activities, and opportunities for industry engagement. Proposed program-level changes are processed through the relevant departmental and faculty governance bodies, with formal approval obtained from the authorised body.

The program maintains supporting records of survey instruments, response summaries, attainment calculations, meeting minutes, and documented improvement actions, as applicable. Subsequent reviews are used to monitor progress and inform future assessment cycles. This process links stakeholder feedback with program review and CQI while distinguishing documented assessment results from improvement actions that are proposed, approved, or implemented.

\section{Design Projects Sponsoring / Supervision Jointly by Faculty and Industry Professionals}\label{sec:c9-design-projects}

The Software Engineering Department at SEECS encourages industry engagement in Final Year Design Projects (FYDPs) to expose students to practical challenges and contemporary industry requirements. Where opportunities arise, industry professionals may contribute by suggesting real-world problem statements, providing technical advice, mentoring student teams, or offering feedback during project development. Faculty supervisors remain responsible for academic supervision, assessment, and ensuring alignment with the programme's learning outcomes. The extent of industry involvement varies across projects and depends on the nature of the project and the arrangements with the relevant organisation. Such engagement can help students strengthen their technical, problem-solving, teamwork, and professional skills while promoting the development of practical solutions to real-world problems. Evidence of industry-sponsored projects, formal joint supervision, or industry participation in project evaluation should be reported where supported by project records. Sample project list is attached as Annexure~\ref{annex:industryproj}.

\section{Industry Collaborations}\label{sec:c9-industry-collab}

Moreover, SEECS and SE Department are also involved in various projects with the industry. The details of some of these projects are given in Table~\ref{tab:colab}, rest are available in evidence folders.

%***************

\begin{longtable}{>{\raggedright\arraybackslash}p{0.04\textwidth} >{\raggedright\arraybackslash}p{0.18\textwidth} >{\raggedright\arraybackslash}p{0.20\textwidth} >{\raggedright\arraybackslash}p{0.50\textwidth}}
\caption{Faculty-Industry Collaborations}\\
\label{tab:colab}\\
\toprule
\textbf{S No} & \textbf{Company} & \textbf{Principal Investigator} & \textbf{Area of Collaboration} \\
\midrule
\endfirsthead
\toprule
\textbf{S No} & \textbf{Company} & \textbf{Principal Investigator} & \textbf{Area of Collaboration} \\
\midrule
\endhead
\midrule
\multicolumn{4}{r}{\emph{Continued on next page}}\\
\endfoot
\bottomrule
\endlastfoot
1 & Walee Technologies (Pvt) Ltd. & Dr. Rafia Mumtaz & Publish Research on business use cases associated with Internet Data Streams \\
%2 & Startify FunPulse Technology (Pvt) Ltd. & Dr. Sohail Iqbal & Efficacy of the new system using quantum safe crypto algorithms to safeguard the intellectual property through novel means on Unity 3D based games development engine \\
2 & Islamabad Diagnostic Centre & Dr. Rafia Mumtaz & The department of R\&D IDC shall provide samples for research and academic activities. \\
3 & OneScreen System & Dr. Aimal Tariq Rextin & Lab + Software Development + Mutual Collaborations \\
%5 & Securiti AI & Dr. Hasan Tahir Butt & IT in Security and Cloud Computing + multiple recruitments \\
%6 & Software Productivity Strategists (SPS) & Dr. Shams Qazi & R\&D Propects + Academia-Industry exchanges for mutual enhancement of skillset and skilled HR \\
4& FFC & Dr. Momina Moetesum & Agri-tech Solutions\\
5 &Wacom & Dr. Faisal Shafait & IoT Paper Device Evaluation and testing on Low resource languages\\
\end{longtable}

%%*************

\section{Research/ Consultancy Projects}\label{sec:c9-consultancy}

SEECS and SE faculty is also actively participating in national/international grants and providing consultancy services to the industry and various other organizations in Pakistan and abroad. Few of the projects by SEECS faculty are mentioned in Table~\ref{tab:consult}.

\begin{longtable}{>{\raggedright\arraybackslash}p{0.05\textwidth} >{\raggedright\arraybackslash}p{0.48\textwidth} >{\raggedright\arraybackslash}p{0.24\textwidth} >{\raggedright\arraybackslash}p{0.17\textwidth}}
\caption{List of Projects / Consultancies}\\
\label{tab:consult}\\
\toprule
\textbf{S No} & \textbf{Project Title} & \textbf{Faculty Name} & \textbf{Funding Agency} \\
\midrule
\endfirsthead
\toprule
\textbf{S No} & \textbf{Project Title} & \textbf{Faculty Name} & \textbf{Funding Agency} \\
\midrule
\endhead
\midrule
\multicolumn{4}{r}{\emph{Continued on next page}}\\
\endfoot
\bottomrule
\endlastfoot
1 & Smart Greenhouse Technologies for Sustainable Agriculture: An Industrial Innovation Partnership between Turkiye and Pakistan & Dr. Momina Moetesum & PSF TUBITAK \\
%2 & Phase - 2, NUST AI Chatbot & Dr. Sohail Iqbal & NUST \\
2 & Early Detection of Forest Health Decline using Remote Sensing Data & Dr. Faisal Shafait & DAAD \\
%4 & CoPEST : Cotton Pest Prevalence Estimation System for Early Warning using UAV \& AIoT & Dr. Muhammad Moazam Fraz & DAAD \\
%5 & An AI enabled quasi-real-time water quality monitoring for early chemical and/or bio-contamination detection & Dr. Muhammad Moazam Fraz & Asi@Connect (South Korea) \\
%6 & ICESCO Chair of Data Science and Analytics for Business & Dr. Muhammad Moazam Fraz - SEECS Dr. Nabeel Anwar – SMME & The Islamic World Educational, Scientific and Cultural Organization (ICESCO), Morocco \\
%7 & XAIBots-Responsibiltiy and Explainability Interactive Bots & Dr. Junaid Younus & BlueDiamond Engineering Austria \\
3 & NeuReader: Eye Tracking Enabled Explainable-AI for Empowering Resource Scarce Neurological Healthcare in Pakistan & Dr. Faisal Shafait & Aston University \\
%9 & SLUM-i: A Remote Sensing Based Intelligent Framework for Detection and Predictive Growth Analysis of Slums for Sustainable Socio-Economic Development of Pakistan & Dr. Muhammad Imran Malik & DAAD \\
%10 & Super-efficient Sustainable Cooling Solution for All Applications (S2Cool) & Dr. Muhammad Moazam Fraz & UKRI \\
%11 & Ba-Ikhtiyar Jawan: Upscaling \& Digitisation of Vocational Education Curriculum in Pakistan & Dr. Muhammad Moazam Fraz & DAAD \\
%12 & CASREW : Uncovering Abiotic Stress in Rain-Fed Wheat & Dr. Muhammad Moazam Fraz & DAAD \\
%13 & CAREON - Continuous AI-Remote E-Health Operations Network & Dr. Nazia Perwaiz & PMU, Saudi Arabia \\
4 & Walee-NUST Center of Excellence for Advanced Social and Digital Media Technologies (Phase 2) & Dr. Rafia Mumtaz & Walee Tech \\
5 & NUST-OneScreen Research Lab & Dr. Aimal Tariq Rextin & One Screen \\
%16 & Open-Source Unified Threat Assessment Suite (UTAS) for Isolated Active Directory (AD) Network & Dr. Hasan Tahir Butt & NESCOM RAC-XI \\
6 & LLM Powered Document Generator and Code Verifier & Dr. Sidra Sultana & NESCOM \\
%18 & Deep Image Matching for Vision-Based 3D Localization of UAVs in GPS-Denied Environments & Dr. Muhammad Moazam Fraz & NESCOM \\
%19 & Characterization of Code Obfuscation and Encryption in Modern Malware & Dr. Mehdi Hussain & NESCOM \\
%20 & FPGA based Emulation of Quantum Computer & Dr. Madiha Khalid & NESCOM \\
%21 & Improving Bilingual Data Quality in Citizen Records \& Semantic Information Retrieval & Dr. Naima Iltaf & NADRA \\
%22 & CompSafeNano Nanoinformation Approaches for Safe-by-Design NanoMaterials & Dr. Muhammad Moazam Fraz & Marie Sklodowaska-Cruie RISE, Harizon 2020 \\
%23 & Towards on-device Classification of Fruit Diseases & Dr. Hasan Ali Khattak & NUST  \\
%24 & Enabling Smart Yield Boosting: Rice Phenology Estimation Using Airborne Multispectral Remote Sensing Data & Dr. Muhammad Moazam Fraz & DAAD  \\
%25 & Vision based Obstacle Avoidance system for Arial Vehicles & Dr. Muhammad Moazma Fraz & NESCOM  \\
%26 & \#15756, Generating Plausible Counterfactual Explanations using Transformers in Natural Language based Financial Forecasting (NLFF) & Dr. Seemab Latif & HEC - NRPU  \\
%27 & Development of Prototype for AI-based Detection of Malware & Dr. Sana Qadir & NUST  \\
%28 & Content Extraction \& Filtering from Video Stream & Dr. Rabia Irfan & NESCOM RAC-IX  \\
7 & Air Quality Assessment based on Industry 4.0 & Dr. Rafia Mumtaz & NUST \\
30 & Transforming the Judiciary System in Pakistan using Artificial Intelligence (Flagship proj) & Dr. Faisal Shafait & NUST \\
%31 & Automatic Classification of Marine Vessels using Acoustic Data & Dr. Qaiser Riaz & NESCOM \\
\end{longtable}

%**************

